\begin{helper}[Product Model Expectation Bound]
Fix \(n\in\mathbb N\), \(I\subseteq\mathrm{Fin}(n)\), parameters
\(\varepsilon,\varepsilon_1,\varepsilon_2\in\mathbb R\), and
\(n_0\in\mathbb N\).  Assume
\[
\noderef{ParameterHierarchyEventualComparisons}
  (\varepsilon,\varepsilon_1,\varepsilon_2,n_0),\qquad n_0<n,
\]
and
\[
|I|=\noderef{TwoBiteNaturalIndependenceScale}(1+\varepsilon,n).
\]
Let
\[
M=\noderef{TwoBiteNaturalReducedVertexCount}(n),\qquad
p_0=\noderef{TwoBiteEdgeProbability}\!\left(\frac12,n\right),
\]
\[
L=\log n,\qquad
C=\noderef{RealChooseTwo}
  (\noderef{TwoBiteNaturalIndependenceScale}(1+\varepsilon,n)).
\]
For every embedding
\[
\pi:\mathrm{Fin}(n)\hookrightarrow\mathrm{Fin}(M)\times\mathrm{Fin}(M),
\]
put
\[
P_R=\{\pi(u)_1:u\in I\},\qquad P_B=\{\pi(u)_2:u\in I\},
\]
\[
q=\noderef{FixedSetProductModelMassFn}(p_0,P_R,P_B),
\qquad
Z_{\mathrm{prod}}=\noderef{FixedSetProductModelVar}(I,\varepsilon,p_0,\pi).
\]
Then the product expectation satisfies
\[
\sum_v\left(\prod_i q(v_i)\right)Z_{\mathrm{prod}}(v)
\le \frac{C}{2L},
\]
where \(v\) ranges over assignments from
\(\mathrm{Fin}(2M)\) to pairs of subsets of \(\mathrm{Fin}(M)\).
\end{helper}
\begin{proof}
The product law uses independent coordinates with one-coordinate mass \(q\).
The hierarchy gives \(0\le p_0\le1\), so
\(\noderef{FixedSetProductModelMass}\) applies to \(P_R\) and \(P_B\).  Thus
\[
q(a)\ge0\quad\text{for every }a,\qquad \sum_a q(a)=1.
\]
Consequently the product weights
\[
W(v)=\prod_{i\in\mathrm{Fin}(2M)}q(v_i)
\]
are nonnegative and have total mass \(1\): summing over the last coordinate
uses \(\sum_aq(a)=1\), then summing over the previous coordinate uses the same
identity again, and iterating over the finite coordinate set leaves \(1\).
Therefore the displayed finite sum is the expectation of
\(\noderef{FixedSetProductModelVar}\) under this product distribution.

We record the marginal calculation used below.  Fix a red coordinate \(x\) and
two distinct red projection vertices \(r_1,r_2\in P_R\).  The event that the
first component of \(v_x\) contains both \(r_1\) and \(r_2\) has probability
\(p_0^2\).  Indeed, in the sum over \(a=(A_R,A_B)\), all terms with
\(A_R\not\subseteq P_R\) or \(A_B\not\subseteq P_B\) are zero.  On the support,
the \(A_B\)-sum is
\[
\sum_{A_B\subseteq P_B}
  p_0^{|A_B|}(1-p_0)^{|P_B|-|A_B|}=1
\]
by the same binomial calculation as in
\(\noderef{FixedSetProductModelMass}\).  In the \(A_R\)-sum, requiring
\(r_1,r_2\in A_R\) factors out \(p_0^2\), and the remaining elements of
\(P_R\setminus\{r_1,r_2\}\) contribute
\[
\sum_B p_0^{|B|}(1-p_0)^{|P_R|-2-|B|}
=(p_0+(1-p_0))^{|P_R|-2}=1.
\]
Thus the marginal probability is \(p_0^2\).  The identical argument for a blue
coordinate says that the second component contains two specified distinct
vertices of \(P_B\) with probability \(p_0^2\).  Since the product weights
factor over coordinates, all other coordinates sum out to \(1\) when this
one-coordinate marginal is computed.

Unfold \(Z_{\mathrm{prod}}\).  A counted pair is a rank-oriented final pair
\((u,w)\) drawn from a two-element subset of \(I\) whose red projections and
blue projections are both
distinct, and which is supported by at least one external red or blue base
vertex in the small class.  Fix such a pair.  Write
\[
r_1=\pi(u)_1,\qquad r_2=\pi(w)_1,\qquad
b_1=\pi(u)_2,\qquad b_2=\pi(w)_2.
\]
The distinct-projection filter gives \(r_1\ne r_2\) and \(b_1\ne b_2\), and
all four vertices lie in \(P_R\) or \(P_B\) as appropriate.  In the product
sample, red witnesses are read only from the first \(M\) coordinates: a red
coordinate can support this pair only if its first component contains both
\(r_1\) and \(r_2\).  For each fixed red coordinate the one-coordinate
marginal above gives probability \(p_0^2\), and the small-class and final-pair
filters can only decrease the event.  Summing over the \(M\) red coordinates
gives a red contribution at most \(Mp_0^2\).

Blue witnesses are read only from the last \(M\) coordinates.  The identical
argument, using the second component and the vertices \(b_1,b_2\), gives a
blue contribution at most \(Mp_0^2\).  Thus
\(\noderef{FixedSetProductModelPairUnionBound}\) bounds the probability that
this fixed pair contributes anywhere by
\[
2Mp_0^2.
\]

From \(\noderef{ParameterHierarchyEventualComparisons}\), in the range
\(n_0<n\) we have
\[
M\le \frac{n}{L^2}
\]
after viewing \(M\) as a real number.  Also
\[
p_0=\frac12\sqrt{\frac{L}{n}}.
\]
Consequently
\[
2Mp_0^2
\le 2\frac{n}{L^2}\cdot\frac{L}{4n}
=\frac1{2L}.
\]

The fixed-size hypothesis gives
\[
|I|=\noderef{TwoBiteNaturalIndependenceScale}(1+\varepsilon,n).
\]
Thus the number of unordered pairs available in \(I\) is at most
\[
C=\noderef{RealChooseTwo}
  (\noderef{TwoBiteNaturalIndependenceScale}(1+\varepsilon,n)).
\]
The last display is nonnegative, so the bound \(1/(2L)\) is also nonnegative.
By \(\noderef{FixedSetProductModelLinearity}\), summing the contribution bound
\(1/(2L)\) over these at most \(C\) pairs gives
\[
\mathbb E_q Z_{\mathrm{prod}}\le \frac{C}{2L}.
\]
\end{proof}
